%==============================================================================
% PARTIDA 04: CONDUCTORES Y CABLES
% Codigos: OE.5.2.3.1 al OE.5.2.3.7
%==============================================================================
\subsection{OE.5.2.3.1 Conductor de Cobre LSOH 1.5 mm2 (Fase + Neutro Alumbrado)}

\subsubsection{Especificaciones Tecnicas}

Conductor de cobre electrolitico recocido para circuitos de alumbrado C1, C3 y C5. Debe cumplir con NTP 370.252, NTP-IEC 60228 (Clases de conductores) y CNE-U Articulo 500.1.

\begin{itemize}[nosep]
  \item \textbf{Material:} Cobre electrolitico 99.9\% de pureza minima.
  \item \textbf{Clase de cableado:} Clase 2 (cableado compacto circular) segun NTP-IEC 60228.
  \item \textbf{Seccion nominal:} 1.5 mm2 (tolerancia: ±2.5\% segun NTP 370.252).
  \item \textbf{Numero de hilos:} 7 hilos, diametro de cada hilo: 0.52 mm.
  \item \textbf{Diametro exterior del conductor:} 3.5 mm ± 0.2 mm.
  \item \textbf{Resistencia electrica maxima a 20 °C:} 12.1 $\Omega$/km.
  \item \textbf{Aislamiento:} Compuesto termoestable LSOH (Low Smoke Zero Halogen), tipo NH-80.
  \item \textbf{Tension nominal:} 450/750 V.
  \item \textbf{Temperatura maxima de operacion:} 80 °C en servicio continuo, 250 °C en cortocircuito (5 s max).
  \item \textbf{Capacidad de corriente en tuberia empotrada:} 15 A a 30 °C (CNE-U Tabla 5.2, 3 conductores cargados).
  \item \textbf{Colores:} Rojo o marron para fase (F), blanco o celeste para neutro (N).
  \item \textbf{Normas:} NTP 370.252, NTP-IEC 60228, NTP-IEC 60332-1 (no propagacion de llama), NTP-IEC 60754 (libre de halogenos).
\end{itemize}

\paragraph{Caida de tension admisible:}
Segun CNE-U Articulo 500.5, la caida de tension maxima en circuitos derivados es 2.5\%. Para el circuito C1 (alumbrado primer piso) con 5 luminarias LED de 12 W cada una (60 W total):

\begin{equation}
\Delta V = \frac{2 \times L \times I \times \rho}{S} = \frac{2 \times 25 \times 0.27 \times 0.0172}{1.5} = 0.15\% \ll 2.5\% \text{ OK}
\end{equation}

Donde $L$ = 25 m (longitud maxima estimada), $I$ = 0.27 A, $\rho$ = 0.0172 $\Omega$mm2/m, $S$ = 1.5 mm2.

\subsubsection{Requisitos de Materiales}

\begin{longtable}{L{5.0cm} L{2.0cm} L{7.5cm}}
\toprule
\textbf{Material} & \textbf{Cant.} & \textbf{Especificacion tecnica} \\
\midrule
Conductor LSOH 1.5 mm2 rojo/marron (fase) & 130 m & Cobre electrolitico clase 2, 450/750 V, NH-80 \\
Conductor LSOH 1.5 mm2 blanco/celeste (neutro) & 110 m & Cobre electrolitico clase 2, 450/750 V, NH-80 \\
Conector tipo "wago" 3 bornes 1.5-2.5 mm2 & 20 pza | Para empalmes en cajas de derivacion \\
Conector tipo "wago" 5 bornes 1.5-2.5 mm2 & 10 pza | Para derivaciones multiples \\
Cinta aislante vinilica 3M Tempflex 1600 & 2 pza | Clase H (105 °C), 600 V, 19 mm x 20 m \\
\bottomrule
\end{longtable}

\subsubsection{Metodo de Ejecucion}

\begin{enumerate}
  \item \textbf{Preparacion:} Verificar que la tuberia este instalada y aprobada (prueba de paso realizada).
  \item \textbf{Corte:} Medir la longitud requerida entre caja de derivacion y punto de salida, agregando 0.30 m en cada extremo para conexiones.
  \item \textbf{Pelado de aislamiento:} Utilizar pelacables automatico ajustado a 1.5 mm2. Pelar 1.5 cm del extremo. No danar los hilos de cobre.
  \item \textbf{Tendido:}
  \begin{enumerate}
    \item Introducir guia de alambre calibre 14 por la tuberia.
    \item Fijar los conductores a la guia con cinta aislante.
    \item Halar suavemente la guia desde el extremo opuesto, guiando los conductores para evitar enredos.
    \item No forzar los conductores; si hay resistencia excesiva, revisar la canalizacion.
  \end{enumerate}
  \item \textbf{Conexiones en cajas de derivacion:}
  \begin{enumerate}
    \item Identificar fase (rojo) y neutro (blanco) en ambos extremos.
    \item Realizar empalmes utilizando conectores tipo "wago" (no se permite cinta aislante como unico medio de empalme en cajas accesibles).
    \item No realizar empalmes dentro de tuberias (CNE-U Articulo 350.10).
  \end{enumerate}
  \item \textbf{Identificacion:} Marcar los conductores en cada caja de derivacion con el numero de circuito (C1, C3, C5).
\end{enumerate}

\subsubsection{Control de Calidad}

\begin{longtable}{L{4.5cm} L{4.5cm} L{4.5cm}}
\toprule
\textbf{Parametro} & \textbf{Metodo de verificacion} & \textbf{Criterio de aceptacion} \\
\midrule
Seccion & Calibrador Vernier sobre cobre pelado & 1.5 mm2 ± 2.5\% (diametro ≈ 1.38 mm) \\
Continuidad & Multimetro digital en modo continuidad & < 1 $\Omega$ por tramo (menos de 100 m) \\
Resistencia de aislamiento & Megohmetro 1000 V entre F-N, F-T, N-T & > 1 M$\Omega$ (recomendado > 20 M$\Omega$) \\
Resistencia electrica & Puente de Wheatley o microhmmetro & < 12.1 $\Omega$/km a 20 °C \\
Identificacion de colores & Inspeccion visual & Rojo/marron = fase, blanco/celeste = neutro \\
\bottomrule
\end{longtable}

\subsubsection{Metodo de Medicion}

Se medira por metro lineal (m) de conductor instalado, medido sobre el eje de la tuberia (longitud desarrollada), incluyendo conexiones y pruebas.

\subsubsection{Unidad de Medida}

m (metro lineal).

\subsubsection{Forma de Pago}

Pago por metro lineal instalado, verificado mediante prueba de continuidad y aislamiento, aprobado por supervision.

%------------------------------------------------------------------------------
\subsection{OE.5.2.3.3 Conductor de Cobre LSOH 2.5 mm2 (Fase + Neutro Tomacorrientes)}

\subsubsection{Especificaciones Tecnicas}

Conductor de cobre LSOH 2.5 mm2 para circuitos de tomacorrientes C2, C4 y C6. La seccion de 2.5 mm2 es la minima exigida por CNE-U para circuitos de tomacorrientes residenciales.

\begin{itemize}[nosep]
  \item \textbf{Seccion nominal:} 2.5 mm2.
  \item \textbf{Numero de hilos:} 7 hilos, diametro cada hilo: 0.67 mm.
  \item \textbf{Diametro exterior:} 4.2 mm ± 0.2 mm.
  \item \textbf{Resistencia electrica maxima:} 7.41 $\Omega$/km a 20 °C.
  \item \textbf{Capacidad de corriente:} 20 A en tuberia empotrada (CNE-U Tabla 5.2).
  \item \textbf{Colores:} Rojo para fase (F), blanco para neutro (N), verde/verde-amarillo para proteccion a tierra (PE).
\end{itemize}

\paragraph{Caida de tension para tomacorrientes:}
Para el circuito C4 (tomacorrientes segundo piso) con mayor carga: 1386 W demandados, 6.30 A, longitud maxima 30 m:

\begin{equation}
\Delta V = \frac{2 \times 30 \times 6.30 \times 0.0172}{2.5} = 2.60\%
\end{equation}

Supera ligeramente el 2.5\% recomendado. Se recomienda verificar en obra y considerar incrementar seccion a 4 mm2 si la longitud real supera los 30 m.

\subsubsection{Requisitos de Materiales}

\begin{longtable}{L{5.0cm} L{2.0cm} L{7.5cm}}
\toprule
\textbf{Material} & \textbf{Cant.} & \textbf{Especificacion tecnica} \\
\midrule
Conductor LSOH 2.5 mm2 rojo (fase) & 320 m & Cobre electrolitico clase 2, 450/750 V \\
Conductor LSOH 2.5 mm2 blanco (neutro) & 130 m & Cobre electrolitico clase 2, 450/750 V \\
Conductor LSOH 2.5 mm2 verde (tierra PE) & 150 m & Cobre electrolitico clase 2, 450/750 V \\
Conector "wago" 3 bornes 2.5 mm2 & 30 pza | Para empalmes en cajas \\
Conector "wago" 5 bornes 2.5 mm2 & 15 pza | Para derivaciones multiples \\
\bottomrule
\end{longtable}

\subsubsection{Metodo de Ejecucion}

Identico al procedimiento de la partida OE.5.2.3.1, con las siguientes particularidades:
\begin{itemize}
  \item Los conductores de 2.5 mm2 se tienden por tuberia PVC pesada (PV-P) de 3/4''.
  \item El conductor verde/verde-amarillo de proteccion a tierra se conecta al borne de tierra de cada tomacorriente y a la barra de tierra del tablero.
  \item No se permite empalmar el conductor de tierra (debe ser continuo desde el tomacorriente hasta la barra de tierra).
\end{itemize}

\subsubsection{Control de Calidad}

\begin{longtable}{L{4.5cm} L{4.5cm} L{4.5cm}}
\toprule
\textbf{Parametro} & \textbf{Metodo de verificacion} & \textbf{Criterio de aceptacion} \\
\midrule
Seccion & Calibrador Vernier & 2.5 mm2 ± 2.5\% (diametro ≈ 1.78 mm) \\
Continuidad & Multimetro & < 1 $\Omega$ por tramo \\
Aislamiento & Megohmetro 1000 V & > 1 M$\Omega$ \\
Polaridad en tomacorrientes & Comprobador de tomacorrientes & Fase en borne L, neutro en N, tierra conectada \\
\bottomrule
\end{longtable}

\subsubsection{Metodo de Medicion}

Se medira por metro lineal (m) de conductor de 2.5 mm2 instalado.

\subsubsection{Unidad de Medida}

m (metro lineal).

%------------------------------------------------------------------------------
\subsection{OE.5.2.3.6 Conductor de Cobre LSOH 10 mm2 (Alimentador General)}

\subsubsection{Especificaciones Tecnicas}

Conductor de cobre LSOH 10 mm2 para el alimentador general desde caja portamedidor hasta TG-01.

\begin{itemize}[nosep]
  \item \textbf{Seccion nominal:} 10 mm2.
  \item \textbf{Numero de hilos:} 7 hilos, diametro cada hilo: 1.35 mm.
  \item \textbf{Diametro exterior:} 6.5 mm ± 0.3 mm.
  \item \textbf{Resistencia electrica maxima:} 1.83 $\Omega$/km a 20 °C.
  \item \textbf{Capacidad de corriente:} 52 A en tuberia empotrada (CNE-U Tabla 5.2).
  \item \textbf{Colores:} Negro para fase, blanco para neutro.
\end{itemize}

\paragraph{Verificacion por caida de tension:}
Demanda maxima: 21.73 A, alimentador: 10 mm2, longitud: 15 m desde medidor hasta TG-01.

\begin{equation}
\Delta V = \frac{2 \times 15 \times 21.73 \times 0.0172}{10} = 1.12\% < 2.5\% \text{ OK}
\end{equation}

\subsubsection{Metodo de Ejecucion}

\begin{enumerate}
  \item Tender conductores por tuberia PVC pesada 1'' desde caja portamedidor hasta TG-01.
  \item Dejar 0.50 m de conductor en cada extremo.
  \item Pelar 2.5 cm de aislamiento para conexion en bornes del interruptor general.
  \item Utilizar terminales de compression tipo ojo para 10 mm2, crimpeando con pinza de compresion hidraulica.
  \item Torque de apriete en bornes: 15 Nm.
\end{enumerate}

\subsubsection{Control de Calidad}

\begin{itemize}
  \item Verificar seccion con calibrador (diametro ≈ 3.57 mm para 10 mm2).
  \item Medir resistencia de aislamiento (megohmetro 1000 V): > 1 M$\Omega$.
  \item Verificar continuidad: < 1 $\Omega$.
  \item Verificar torque de apriete de conexiones (15 Nm).
\end{itemize}

\subsubsection{Metodo de Medicion}

Se medira por metro lineal (m) de conductor de 10 mm2 instalado.

\subsubsection{Unidad de Medida}

m (metro lineal).

%------------------------------------------------------------------------------
\subsection{OE.5.2.3.7 Conductor de Cobre Temple Blando 6 mm2 (Conexion Puesta a Tierra)}

\subsubsection{Especificaciones Tecnicas}

Conductor de cobre electrolitico de temple blando, seccion 6 mm2, desnudo (sin aislamiento), para conexion del electrodo de puesta a tierra a la barra de tierra del TG-01. Debe cumplir con CNE-U Articulo 600.2 (Conductores de puesta a tierra).

\begin{itemize}[nosep]
  \item \textbf{Seccion:} 6 mm2 (minimo exigido por CNE-U para conductores de proteccion residencial).
  \item \textbf{Temple:} Blando (recocido), para facilitar el tendido.
  \item \textbf{Resistencia electrica:} 3.08 $\Omega$/km a 20 °C.
  \item \textbf{Resistencia a la corrosion:} El cobre desnudo es resistente a la corrosion en suelos de pH 5.5 a 8.5. Para suelos agresivos, utilizar conductor con aislamiento PVC.
\end{itemize}

\subsubsection{Metodo de Ejecucion}

\begin{enumerate}
  \item Tender el conductor de cobre desnudo de 6 mm2 desde la barra de tierra del TG-01 hasta el electrodo de puesta a tierra.
  \item Proteger el conductor con tuberia PVC de 3/4'' en tramos expuestos (sobre piso, en muros).
  \item Conectar a la barra de tierra del TG-01 mediante terminal de compression tipo ojo para 6 mm2.
  \item Conectar al electrodo (varilla de cobre) mediante conector split-bolt de bronce, aplicando gel conductor.
  \item No se permiten empalmes en el conductor de tierra (CNE-U Articulo 600.8).
\end{enumerate}

\subsubsection{Control de Calidad}

\begin{itemize}
  \item Verificar continuidad electrica entre barra de tierra y electrodo (ideal: < 0.5 $\Omega$).
  \item Medir resistencia de puesta a tierra con telurómetro: < 15 $\Omega$.
  \item Inspeccionar la proteccion mecanica en tramos expuestos.
\end{itemize}

\subsubsection{Metodo de Medicion}

Se medira por metro lineal (m) de conductor de 6 mm2 instalado, incluyendo conectores y proteccion.

\subsubsection{Unidad de Medida}

m (metro lineal).
